Galatea és el quart satèl·lit de Neptú més allunyat del planeta. Va ser descobert per la sonda espacial \textit{Voyager 2} l'any 1989. Suposem que l'òrbita que descriu és circular.

\figura[0.35]{fig1.png}
{\centering Imatge de Galatea vista per la sonda \textit{Voyager 2}\par}

\begin{apartat}{1}
Calculeu la velocitat lineal orbital de Galatea en el sistema de referència centrat en Neptú i calculeu la massa de Neptú.
\end{apartat}

\begin{apartat}{1}
Calculeu el valor de la intensitat de camp gravitatori que Neptú crea a la seva pròpia superfície.
\end{apartat}

\dades{%
  Període de l'òrbita de Galatea, $T_\mathrm{Galatea} = 0,428$ dies\\
  Radi de l'òrbita de Galatea, $R_\mathrm{Galatea} = 6,20\times 10^{4}\ \mathrm{km}$\\
  Radi de Neptú, $R_{\text{Neptú}} = 2,46\times 10^{4}\ \mathrm{km}$\\
  $G = 6,67\times 10^{-11}\ \mathrm{N\,m^2\,kg^{-2}}$}
